\documentclass[11pt]{article}

\usepackage[T1]{fontenc}
\usepackage[utf8]{inputenc}
\usepackage{amsmath,amssymb,amsthm}
\usepackage{mathtools}
\usepackage[a4paper,margin=1in]{geometry}
\usepackage{hyperref}

\newcommand{\E}{\mathcal E}
\newcommand{\Var}{\mathsf{Var}}
\newcommand{\B}{\mathsf{B}}
\newcommand{\Tel}{\mathsf{T}}
\newcommand{\Ty}{\mathsf{Ar}}
\newcommand{\arity}{\mathsf{ar}}
\newcommand{\coh}{\mathsf{coh}}
\newcommand{\ext}{\mathbin{\triangleleft}}
\newcommand{\into}{\longrightarrow}
\newcommand{\id}{\mathrm{id}}

\title{The Mathematical Content of \texttt{diamondAt}}
\author{}
\date{}

\begin{document}
\maketitle

\section*{The point}

The theorem \texttt{diamondAt} is an interchange law for hereditary
substitution.  It says that the following two operations commute:

\begin{enumerate}
\item replace variables from one context \(D\) by expressions over \(D'\);
\item replace variables from another context \(S\) by expressions over \(S'\).
\end{enumerate}

There is one important subtlety.  The fillers used to replace the \(S\)-variables
may themselves contain variables from \(D\).  Therefore, if we first substitute
\(D\)-variables, then the \(S\)-fillers must also be transformed by that same
substitution.  The theorem says that this is the only correction needed.

\section*{Notation}

We write \(\E(\Gamma)\) for expressions in context \(\Gamma\).  If \(\alpha\) is
an arity, then
\[
  \Gamma \ext \alpha
\]
means the context \(\Gamma\) extended by one binder of arity \(\alpha\).  This is
written \(\Gamma \mathbin{\texttt{::}} \alpha\) in Lean.

The operation \(\bowtie\) appends telescopes.  We use the following names for
the contexts appearing in the theorem:
\[
  P = pre,\qquad
  D = dom,\qquad
  D' = cod,\qquad
  T = \tau,\qquad
  S = src,\qquad
  S' = dst.
\]
The list \(\upsilon\) of already-opened binders gives a telescope
\[
  U = \Tel(\upsilon).
\]
In Lean, \(U\) is \(\mathsf{Tele.ofList}\,\upsilon\).

The expression to which the theorem applies lives in the context
\[
  ((((P \bowtie D) \bowtie T) \bowtie S) \bowtie U).
\]
Thus it may mention variables from five regions:
\[
  P,\qquad D,\qquad T,\qquad S,\qquad U.
\]

\section*{The two substitutions}

The first substitution is
\[
  \sigma : D \leadsto P \bowtie D'.
\]
This means that for every variable \(z : D \ni \beta\), we have a filler
\[
  \sigma(z) \in \E((P \bowtie D') \ext \beta).
\]
It replaces \(D\)-variables while preserving the prefix \(P\).

The second substitution is
\[
  \kappa : S \leadsto ((P \bowtie D) \bowtie T) \bowtie S'.
\]
Thus for every \(x : S \ni \beta\),
\[
  \kappa(x)
  \in
  \E(((((P \bowtie D) \bowtie T) \bowtie S') \ext \beta)).
\]
This replaces \(S\)-variables by expressions that still live over \(D\).  So
\(\kappa\) is a substitution before the change of target context induced by
\(\sigma\).

\section*{Action at a passive depth}

If \(A\) is a passive telescope, then \(\sigma\) acts at depth \(A\) as a map
\[
  \sigma_A :
  \E((P \bowtie D) \bowtie A)
  \into
  \E((P \bowtie D') \bowtie A).
\]
The word ``passive'' means that variables from \(A\) are not substituted.  They
are merely carried along.  Only variables from \(D\) are active for \(\sigma\).

Similarly, \(\kappa\) acts at passive depth \(U\) as a map
\[
  \kappa_U :
  \E(((((P \bowtie D) \bowtie T) \bowtie S) \bowtie U))
  \into
  \E(((((P \bowtie D) \bowtie T) \bowtie S') \bowtie U)).
\]
Only variables from \(S\) are active for \(\kappa\); everything else is passive.

\section*{The pushforward of a substitution}

Since the fillers of \(\kappa\) may contain \(D\)-variables, applying \(\sigma\)
to an expression also transports the substitution \(\kappa\).  We write this
pushforward as
\[
  \sigma_\ast\kappa : S \leadsto ((P \bowtie D') \bowtie T) \bowtie S'.
\]
It is defined pointwise by
\[
  (\sigma_\ast\kappa)(x)
  =
  \sigma_{T \bowtie S' \ext \beta}(\kappa(x)),
  \qquad x : S \ni \beta.
\]
The passive depth is \(T \bowtie S' \ext \beta\), because the filler
\(\kappa(x)\) is an expression over
\[
  ((P \bowtie D) \bowtie (T \bowtie S' \ext \beta)).
\]
After \(\sigma\) acts, this becomes an expression over
\[
  ((P \bowtie D') \bowtie (T \bowtie S' \ext \beta)),
\]
which is exactly the required type of a filler for \(x\) in the transformed
substitution.

\section*{The theorem as an equation}

Let
\[
  e \in \E(((((P \bowtie D) \bowtie T) \bowtie S) \bowtie U)).
\]
Then \texttt{diamondAt} says:
\[
\boxed{
  (\sigma_\ast\kappa)_U
  \bigl(
    \sigma_{(T \bowtie S) \bowtie U}(e)
  \bigr)
  =
  \sigma_{(T \bowtie S') \bowtie U}
  \bigl(
    \kappa_U(e)
  \bigr).
}
\]

This is the clean mathematical statement.  The left-hand side first applies
\(\sigma\) to \(e\), then substitutes the remaining \(S\)-variables using the
pushforward \(\sigma_\ast\kappa\).  The right-hand side first substitutes the
\(S\)-variables using \(\kappa\), then applies \(\sigma\).

\section*{The same equation as a square}

The theorem says that the following square commutes:
\[
\begin{array}{ccc}
\E(((((P \bowtie D) \bowtie T) \bowtie S) \bowtie U))
& \xrightarrow{\ \sigma_{(T \bowtie S) \bowtie U}\ } &
\E(((((P \bowtie D') \bowtie T) \bowtie S) \bowtie U))
\\[0.8em]
\downarrow{\kappa_U} && \downarrow{(\sigma_\ast\kappa)_U}
\\[0.8em]
\E(((((P \bowtie D) \bowtie T) \bowtie S') \bowtie U))
& \xrightarrow{\ \sigma_{(T \bowtie S') \bowtie U}\ } &
\E(((((P \bowtie D') \bowtie T) \bowtie S') \bowtie U)).
\end{array}
\]

The top arrow changes \(D\) into \(D'\).  The left arrow changes \(S\) into
\(S'\).  The bottom arrow changes \(D\) into \(D'\) after \(S\) has already
been changed into \(S'\).  The right arrow changes \(S\) into \(S'\) after the
fillers for \(S\) have themselves been pushed forward along \(\sigma\).

\section*{Why the theorem is not trivial}

If the fillers of \(\kappa\) did not contain variables from \(D\), the theorem
would be close to ordinary commutation of independent substitutions.  But here
the fillers of \(\kappa\) are expressions over
\[
  ((P \bowtie D) \bowtie T) \bowtie S',
\]
so they may contain \(D\)-variables.  Therefore substituting \(D\)-variables
before substituting \(S\)-variables changes the \(S\)-fillers themselves.

The pushforward \(\sigma_\ast\kappa\) is precisely the bookkeeping that accounts
for this dependency.  The theorem says:
\[
  \text{the dependency of \(\kappa\) on \(D\) is handled pointwise by \(\sigma\).}
\]

\section*{The five head cases}

The context of \(e\) decomposes into five regions:
\[
  ((((P \bowtie D) \bowtie T) \bowtie S) \bowtie U).
\]
Consequently, the head variable of an application in \(e\) belongs to exactly
one of:
\[
  U,\qquad S,\qquad T,\qquad D,\qquad P.
\]
The proof of the theorem is a structural recursion on \(e\), with a case split
on these five possibilities.

\paragraph{Head in \(U\).}
This is a local binder introduced by recursive descent.  Neither substitution
acts on it.  Both sides rebuild the same head and use the induction hypothesis
on all arguments.

\paragraph{Head in \(S\).}
This is an active variable for \(\kappa\).  On the right side, \(\kappa\) fires
immediately.  On the left side, \(\sigma\) first acts on the ambient expression,
and then the pushforward \(\sigma_\ast\kappa\) fires.  The two results are equal
because \(\sigma_\ast\kappa\) was defined by applying \(\sigma\) to the fillers
of \(\kappa\).  This branch recursively invokes the same theorem on the
selected filler \(\kappa(x)\).

\paragraph{Head in \(T\).}
This head is passive for both substitutions.  Both sides preserve it and
recurse on the arguments.

\paragraph{Head in \(D\).}
This is an active variable for \(\sigma\).  Firing \(\sigma\) creates a
one-binder instantiation from the arguments of the application.  This is the
only branch where the companion theorem for one-binder instantiation is needed.

\paragraph{Head in \(P\).}
This is part of the untouched prefix.  Both substitutions preserve it and
recurse on the arguments.

\section*{Why the formal statement has coherence data}

The mathematical equation above suppresses associativity bookkeeping for
\(\bowtie\).  Lean cannot suppress it completely.  The formal theorem therefore
carries witnesses saying that the decompositions
\[
  T \bowtie S
  \qquad\text{and}\qquad
  T \bowtie S'
\]
are proper and that injections into larger contexts compute with the intended
parenthesization.  These witnesses do not change the mathematical content of
the theorem.  They only ensure that the formal proof can recognize, for example,
that a head from \(S\) remains a head from \(S\) after it is embedded through
\[
  (P \bowtie D) \bowtie (T \bowtie S).
\]

\section*{Role in substitution composition}

The public composition theorem says that acting by a composite substitution is
the same as acting by one substitution and then the other.  The difficult case
is when the first substitution fires on an application head.  Firing creates a
fresh one-binder instantiation, and the remaining substitution must commute
past that instantiation.

The theorem described here is the general interchange principle that makes that
step possible.  In short:
\[
  \text{hereditary substitution respects substitution composition.}
\]

\end{document}
